\chapter{Modo caza}
\label{sec:man-caza}

\section{Ejecutar un hunt}

Desde el panel \emph{Hunt}:

\begin{enumerate}
    \item Escribir la hipótesis en DSL.
    \item Opcional: \code{min_score} (default 0.0) — descarta paths
    por debajo de ese score.
    \item Opcional: \code{max_results} (default 500) — corta la
    salida.
    \item \emph{Run}.
\end{enumerate}

La respuesta es una tabla de paths, cada uno con:

\begin{itemize}
    \item El score compuesto en $[0, 1]$.
    \item El desglose de los 5 componentes.
    \item Los nodos del path, en orden.
    \item La ventana temporal (primer y último evento).
\end{itemize}

\section{Paginación}

Si hay más resultados que \code{max_results}, la UI pagina por 50.
Podés cambiar el tamaño. Los resultados se persisten en la sesión;
podés volver a ellos con \emph{Historial de hunts}.

\section{Scoring: los 5 componentes}

El score compuesto es una combinación ponderada de:

\begin{enumerate}
    \item \textbf{Entity rarity} — qué tan infrecuentes son los
    nodos individualmente. Valores altos: entidades poco vistas.
    \item \textbf{Edge rarity} — qué tan infrecuentes son las
    transiciones. Valores altos: secuencia inusual.
    \item \textbf{Neighborhood concentration} — qué tan uniforme es
    el vecindario de cada nodo. Valores altos \emph{penalizan}
    (centralidad estructural, no anomalía real).
    \item \textbf{Temporal novelty} — qué tan recientemente
    aparecieron las entidades. Valores altos: aparecieron al final
    de la ventana.
    \item \textbf{GNN threat} — score de modelo GNN si está
    cargado (por defecto desactivado).
\end{enumerate}

Pesos por defecto: 0.25 / 0.30 / 0.25 / 0.20 / 0.00. Podés
ajustarlos desde \emph{Scoring \textrightarrow{} Pesos} o dejarlos
en su default.

\section{Interpretar un path}

Un path con score 0.73 y el siguiente desglose:

\begin{lstlisting}[basicstyle=\ttfamily\small]
entity_rarity:              0.82
edge_rarity:                0.91
neighborhood_concentration: 0.45
temporal_novelty:           0.64
gnn_threat:                 0.00
\end{lstlisting}

Se lee: las entidades son raras (0.82), la secuencia es rara
(0.91), el vecindario no es particularmente concentrado (0.45,
bien), y las entidades aparecieron relativamente tarde (0.64).
Score final alto, consistente con actividad sospechosa genuina.

Contraejemplo útil: un path con scores altos en rarity pero
\emph{también} alto en \code{neighborhood_concentration} suele ser
una entidad hub (DC, proxy, gateway) que aparece rara simplemente
porque conecta muchas cosas. No es anomalía, es topología.

\section{Workflow de priorización}

\begin{enumerate}
    \item Ordenar por score descendente.
    \item Tomar los top-10. Para cada uno:
    \begin{itemize}
        \item Leer el desglose: ¿qué componente domina?
        \item Si \code{neighborhood_concentration} es muy alto,
        probablemente es ruido estructural; marcar y seguir.
        \item Si \code{entity_rarity} y \code{edge_rarity} son
        altos y \code{neighborhood_concentration} es bajo, es
        candidato fuerte. Explorar.
    \end{itemize}
    \item Para candidatos fuertes: doble click $\to$ modo
    explorador (cap.~\ref{sec:man-explorador}).
    \item Anotar hallazgos.
\end{enumerate}

\section{Explicar un score}

Botón \emph{Explain} sobre un path abre la vista de narrativa:
texto humano que describe por qué ese path obtuvo ese score. No
reemplaza al desglose numérico — lo acompaña.

\section{Diferencial entre dos hunts}

\emph{Hunt \textrightarrow{} Diff} permite comparar dos hunts del
historial: qué paths aparecieron, cuáles desaparecieron, cuáles
cambiaron de score. Útil para:

\begin{itemize}
    \item Comparar hunts antes y después de ingestar más datos.
    \item Comparar hipótesis levemente distintas.
    \item Validar que un cambio de mapping no alteró resultados
    críticos.
\end{itemize}

\section{Recomputar scoring}

Si cambias los pesos (o habilitas GNN) sobre resultados ya
existentes, usar \emph{Scoring \textrightarrow{} Recompute}. No
vuelve a correr el hunt: sólo re-evalúa los paths ya devueltos con
los pesos nuevos. Barato.
